\section{Biện pháp thứ hai: bắt tín hiệu lỗi giữ nguyên độ lớn}
\label{sec:ch03-regularizer}

\index{regularizer!giữ norm}%
Clipping chặn phía nổ. Nó không làm gì cho phía tắt, và bài biết điều đó, nên
họ đưa thêm một biện pháp thứ hai nhắm đúng vào phía còn lại. Biện pháp ấy là
một số hạng phạt:

\begin{equation}
  \Omega = \sum_k \Omega_k
         = \sum_k \left(
             \frac{\norm{ \pd{L}{\vect{a}_{k+1}} \jac{\vect{a}_{k+1}}{\vect{a}_k} }}
                  {\norm{ \pd{L}{\vect{a}_{k+1}} }}
             - 1
           \right)^{2}.
  \label{eq:ch03-omega}
\end{equation}

Đọc nó như một tỷ số. Tử số là tín hiệu lỗi sau khi đi lùi thêm một bước, mẫu
số là chính tín hiệu ấy trước bước đó. Tỷ số là thứ mà bước đó làm với độ lớn
của tín hiệu. Bắt tỷ số bằng 1 ở mọi bước là bắt đường đi ngược giữ nguyên
norm. Bình phương làm cho co và giãn đắt ngang nhau.

Tôi in ra chính các tỷ số ấy, trước khi bình phương, để xem số hạng phạt nhìn
thấy gì. Ba mạng, ba \tn{bán kính phổ}{spectral radius}, 32 đơn vị ẩn, 30 bước:

\begin{measured}{tỷ số từng bước mà \(\Omega\) đo}
\begin{minted}[fontsize=\footnotesize]{text}
 radius k=0         k=10        k=28             mean      Omega
  0.500 0.488384    0.500000    0.500000     0.499463   7.516259
  0.900 0.879230    0.897886    0.899925     0.895816   0.326728
  1.000 0.977480    0.982812    0.978343     0.978175   0.014876
\end{minted}

Đo bằng \code{experiments/ch03\_regularizer.py} tại tag \code{ch03}.
\end{measured}

Ở hai bán kính đầu, tỷ số bám sát chính bán kính phổ: trung bình 0.499463 ở 0.5
và 0.895816 ở 0.9. Ở bán kính 1.0 thì trung bình chỉ 0.978175, nằm dưới, và nằm
dưới vì đúng lý do mục~\ref{sec:ch03-tich-jacobian} đã đo: \(\sigma'\) của tanh
nhỏ hơn 1 ở mọi trạng thái khác gốc. \(\Omega\) giảm đơn điệu khi mạng tiến về
chỗ giữ được norm: 7.516259 ở bán kính 0.5, xuống 0.326728 ở 0.9, xuống 0.014876
ở 1.0. Số hạng phạt đo đúng thứ nó tuyên bố đo.

\subsection*{Vì sao chỉ một trong hai còn được dùng}

\index{regularizer!chi phí}%
Câu trả lời ngắn không phải là \enquote{cái kia không chạy}. Nó chạy, và bài có
số để chứng minh. Trên bài toán thứ tự thời gian, với chuỗi dài hơn 20 bước, cả
stochastic gradient descent, viết tắt là SGD, lẫn SGD cộng clipping đều không
giải được; SGD cộng clipping cộng
regularizer giải được với tỷ lệ thành công 100\% tới 200 bước, bằng một model
duy nhất dùng cho mọi độ dài từ 50 tới 200~\autocite{pascanu2013difficulty}. Đó là khoảng cách giữa không làm
được và làm được, chứ không phải một cải thiện vài phần trăm.

Câu trả lời nằm ở giá phải trả, và giá ấy lộ ra ngay khi viết code. Nhìn
lại~\eqref{eq:ch03-omega}: mỗi số hạng cần \(\pd{L}{\vect{a}_{k+1}}\),
tức tín hiệu lỗi tại từng bước. Backward pass thường có tính ra các vector đó,
nhưng tính xong thì vứt đi; giữ lại toàn bộ chúng là một khoản bộ nhớ tỷ lệ với
độ dài chuỗi. Rồi còn phải đẩy đạo hàm của chính \(\Omega\) ngược vào
\(\mat{W}_{hh}\), tức đạo hàm bậc hai qua một đồ thị đã trải dài theo thời gian.

Bài không trả đủ giá đó. Các tác giả lấy đường tắt và ghi rõ trong bài: chỉ dùng
đạo hàm riêng tức thời, coi \(\vect{a}_k\) và \(\pd{L}{\vect{a}_{k+1}}\)
là hằng. Trong repo companion, đường tắt ấy là hai lời gọi \code{detach}:

\begin{minted}{python}
for k in range(len(backward_signals) - 1):
    signal = backward_signals[k + 1].detach()
    denominator = torch.linalg.vector_norm(signal)
    if denominator.item() == 0.0:
        continue
    jacobian = model.jacobian_at(states[k].detach())
    pushed = signal @ jacobian
    ratio = torch.linalg.vector_norm(pushed) / denominator
    terms.append((ratio - 1.0) ** 2)
\end{minted}

Đặt cạnh nhau thì đánh đổi rõ. Clipping là bốn dòng, không cần gì mà backward
pass chưa có, và không có siêu tham số nào ngoài một ngưỡng mà bài bảo là không
nhạy. Regularizer cần giữ tín hiệu lỗi của từng bước, cần một lượt truyền thứ
hai, cần một hệ số \(\alpha\), và trong các thí nghiệm của chính bài thì hệ số
ấy còn phải theo lịch giảm dần theo epoch trên các tác vụ thật.

Một biện pháp đi vào thư viện chuẩn. Biện pháp kia phải viết tay mỗi lần dùng.
Đó là toàn bộ khác biệt, và nó không phải khác biệt về chất lượng khoa học.

Còn một lý do nữa, và nó là lý do sâu hơn: cùng năm ấy, hướng còn lại của bài
toán đã có một lời giải khác, cũ hơn bài này mười sáu năm, và lời giải đó không
sửa thuật toán huấn luyện mà sửa chính kiến trúc. Chương~\ref{ch:lstm} đọc nó.
